\section{The scalar Gaussian-message predecessor}
\label{app:scalar-prior}
This appendix is a reduction to prior work, not an additional novelty claim.
We state exactly which interpretation of \citet{Mahalanabis2012} is used.
Their Section 3.2, Theorem 43 approximates total conditional variance on a
Gaussian precision graph of treewidth $w$, with runtime
\begin{equation}
 \left(\frac{Nw\,\mathrm{cond}(\Lambda)}{\alpha}\right)^{O(w^2)}k^2
       +N2^{O(w^3)}
 \label{eq:gmrf-runtime}
\end{equation}
for relative variance error $1+\alpha$. Here $N$ counts Gaussian vertices,
not candidate times. The theorem as printed sums all unobserved variances
and permits all nodes as observations. We need only a single unobservable
target, and only designated observation candidates.

\subsection{The focused recurrence and a local-cost issue}
Their Definition 18, equation (37), assigns a conditional-variance cost to
the private variables of a subtree. Equations (39)--(40) combine those local
costs with child messages, while (41)--(43) pass separator precision by
observation and Schur complementation. Rounded counterparts are (44)--(48).
Restricting every local observed subset to designated candidates simply
removes choices; the proof's rounding step retains the same observed subset,
so every allowed optimal witness remains allowed. At the level of the covariance costs intended by Definition 18, replacing
each local trace by a nonnegative weighted sum of its private-coordinate variances
preserves relative-error comparisons: if
$e^{-a}M\preceq\widehat M\preceq e^aM$ with $M\succ0$, inversion gives
the same factors for each covariance diagonal and their nonnegative weighted
sum. The addition, observation, marginalization and rounding comparisons in
Observation 33 and Lemmas 41--42 then give the same induction. Zero costs
are compared by inequalities without division. This is an adaptation of the
proof, not the wording of the stated theorem.

One printed local-cost formula needs care. Taking an inverse after extracting
a private principal block generally differs from extracting that block after
inverting the full unobserved bag precision. For example,
\[
 M=\begin{pmatrix}2&1\\1&3\end{pmatrix},\qquad
 (M^{-1})_{11}=3/5\ne (M_{11})^{-1}=1/2.
\]
For the focused specialization we can avoid a nonzero cost in this disputed
form. Retain the target $\vartheta$ in every separator, direct messages
toward a singleton bag $\{\vartheta\}$ followed by an empty root, and give
all earlier private variables weight zero. The only positive cost is then
the target variance in that singleton bag, whose separator is empty and
where the two inverse/block orders coincide. Zero-weight local costs are
set to zero without evaluating an unnecessary inverse. Information messages
still use the stated Schur complements. Binary copies of bags retain
$\vartheta$ until the singleton and preserve this property. Thus the
focused construction does not depend on that erroneous local-cost identity.

\subsection{Rational normalization and regularization}
Take the scalar-observation class of Theorem~\ref{thm:trace-fptas}, with one
parameter ($p=1$) and cardinality as the only constraint. Cardinalities $k<0$
or $k>n$ are infeasible; for $k=0$, return the empty schedule. If all
sensitivities vanish, every feasible schedule is optimal. In the remaining
case, assume $1\le k\le n$ and at least one nonzero sensitivity $f_t$.
Choose positive powers of two $s,d_t$ with
\[
 s^2\le r_{\min}<4s^2,\qquad d_t^2\le r_t<4d_t^2.
\]
Define
\[
 U_t=X_t/s,\quad Z_t=Y_t/d_t,\quad c_t=s/d_t,
 \quad h=\max_t|f_t/d_t|>0,\quad g_t=f_t/(d_th).
\]
Then $|c_t|\le2$, $|g_t|\le1$, some $|g_t|=1$, and the normalized
observation-noise variances lie in $[1,4)$. The normalized parameter is
$\vartheta=h\theta$; its data information equals original data information
divided by $h^2$, which preserves rankings and relative ratios. The latent
covariance $K=\Cov(U)$ satisfies
\[
 |K_{ij}|\le4B_0\rho_0^{|i-j|},\qquad
 \norm K_2\le4B_0(1+\rho_0)/(1-\rho_0).
\]
For zero contraction the same formulas apply directly.

To allow zero original process variances, add independent variance $\zeta$
to each normalized initial/process innovation. Write $\mathsf A$ for the
strict lower shift with entries $a_t$. Its norm is at most $\rho_0$; thus
$T=(I-\mathsf A)^{-1}$ satisfies
\[
 \frac{\zeta}{(1+\rho_0)^2}I\preceq E:=\zeta TT^\top
       \preceq\frac{\zeta}{(1-\rho_0)^2}I.
\]
The lower bound follows from $\norm{T^{-1}}_2\le1+\rho_0$ and the upper
one from the Neumann series. Hence $K_+=K+E\succ0$ has polynomially
controlled spectrum. Let $R'$ be the normalized observation-error covariance
and $C=\diag(c_t)$. Since $R'\succeq I$ and $\norm C_2\le2$,
\[
 R'\preceq R'_+:=R'+CEC^\top
 \preceq R'+\frac{4\zeta}{(1-\rho_0)^2}I.
\]
Choose $\tau=\varepsilon/4$ and
$\zeta=\tau(1-\rho_0)^2/4$. Every selected principal covariance obeys
$R'_{SS}\preceq R'_{+,SS}\preceq(1+\tau)R'_{SS}$, so for
$I'(S)=g_S^\top(R'_{SS})^{-1}g_S$ and its regularized counterpart,
\begin{equation}
 I'(S)/(1+\tau)\le I'_+(S)\le I'(S).
 \label{eq:gmrf-regularization}
\end{equation}
All these changes are rational and have polynomial encoding length.

\subsection{Constant width and conversion of the objective}
Introduce an artificial independent target prior $\vartheta\sim N(0,1)$
and the model
\[
 U\sim N(0,K_+),\qquad Z_t=g_t\vartheta+c_tU_t+\nu_t,
 \qquad 1\le\Var(\nu_t)<4.
\]
Only the $Z_t$ are candidates; only $\vartheta$ has positive variance weight.
Then $v(S)=\Var(\vartheta\mid Z_S)=1/(1+I'_+(S))$. The Gaussian precision
graph has bags $\{\vartheta,U_t,U_{t+1},Z_t\}$ for $t<n$, followed by
$\{\vartheta,U_n,Z_n\}$ and the singleton target/empty root above.
These cover the latent-chain edges and each observation-density triangle,
and satisfy running intersection. Their width is at most three.

For conditioning, start from independent $(\vartheta,U,\nu)$ with covariance
$\diag(1,K_+,D)$, $I\preceq D\prec4I$. The shear to
$(\vartheta,U,Z)$ and its inverse have norms at most
$1+\sqrt{n+4}$, because the added bottom block $[g,C]$ has norm at most
$\sqrt{n+4}$. The joint covariance and inverse precision therefore have
condition number bounded by
\begin{equation}
 (1+\sqrt{n+4})^4
 \frac{\max\{4,\,4B_0(1+\rho_0)/(1-\rho_0)
                       +\zeta/(1-\rho_0)^2\}}
      {\min\{1,\,\zeta/(1+\rho_0)^2\}}
 =O_{\rho_0,B_0}(n^2/\varepsilon).
 \label{eq:gmrf-condition}
\end{equation}
In particular, normalization removes arbitrarily varying original noise and
sensitivity scales from this bound.

At an index with $|g_t|=1$, the regularized observation-error variance is
at most $D_0=4+16B_0+4\zeta/(1-\rho_0^2)$. A best nonempty schedule thus
has $I'_{+,*}\ge1/D_0$: start with that observation and use monotonicity to
pad to cardinality $k$. Apply the focused predecessor with
$\alpha=\varepsilon/[4(D_0+1)]$. Its variance ratio implies
\[
 I'_+(\widehat S)\ge\frac{I'_{+,*}-\alpha}{1+\alpha}
 \ge\frac{1-\alpha D_0}{1+\alpha}I'_{+,*}
 \ge(1-\varepsilon/4)I'_{+,*}.
\]
An at-most-$k$ output can be padded to exactly $k$, since observing more
variables cannot reduce information. Combining with
\eqref{eq:gmrf-regularization} yields
\[
 I'(\widehat S)\ge
 \frac{1-\varepsilon/4}{1+\varepsilon/4}\max_{|S|=k}I'(S)
 \ge(1-\varepsilon)\max_{|S|=k}I'(S).
\]
Multiplication by $h^2$ and addition of a common nonnegative scalar prior
preserve this guarantee. Equations~\eqref{eq:gmrf-runtime} and
\eqref{eq:gmrf-condition} make the predecessor's operation bound polynomial
for this normalized class.

This establishes the mathematical reduction in the predecessor's
spectral-rounding model. It does not silently supply a full rational
implementation of that rounding net. Our explicit bit-complexity result is
proved independently in Theorem~\ref{thm:trace-fptas}; the remaining numerical
model distinction is no basis for claiming a first scalar approximation
scheme. The reduction in this appendix uses cardinality alone and does not
justify padding under spacing or arbitrary side constraints.

\section{Proof of the represented-matroid cover}
\label{app:matroid-proof}
We prove Theorem~\ref{thm:matroid-spectral-set} in full to expose both exact
attainability and representation assumptions. If $q=\rank A=0$, the only
base is empty, and it is returned. Otherwise select $q$ independent rows
of $A$ and discard the redundant rows. This preserves its column matroid.
Apply Lemma~\ref{lem:spectral-normalization}, retaining its exact
range and magnitude tests. For a rank-$r$ trial, let $E'$ be the retained
elements and $F$ the distinct forced owners. Reject the trial unless
$F\subseteq E'$, $A_F$ has independent columns, and $\rank A_{E'}=q$.
The last condition is essential: a base of a lower-rank restriction would
not be a base of the original matroid.

\subsection{Contraction preserving actual bases}
Extend the columns $A_F$ to a square basis matrix $D$ using columns from
$E'$, with $F$ first. Apply $D^{-1}$ to $A_{E'}$, remove the first $|F|$
rows and columns indexed by $F$, and call the result $A_c$. It has full
row rank $q'=q-|F|$. The removed directions are exactly $\range A_F$;
therefore optional columns are independent modulo that span precisely when
their images in $A_c$ are independent. It follows that
\begin{equation}
 B'\text{ is a base of }A_c
 \quad\Longleftrightarrow\quad
 F\cup B'\text{ is an original base contained in }E'.
 \label{eq:matroid-contraction}
\end{equation}
These operations are rational with polynomial bit complexity. Clear
all denominators by a common nonzero integer if an integer representation
is desired. If $q'=0$, return $F$ for this trial. It is already an original
base and satisfies the forced normalization floor.

\subsection{Exact profile detection and witness recovery}
Otherwise put $h=\eta/(rq)$ and round only optional atoms by
$z_{e,ij}=\lfloor A_{e,ij}/h\rfloor$, $i\le j$. Keep the prior and forced
sum exact. There are $d=r(r+1)/2$ signed coordinates. Put
$L_* =\lceil4p/h\rceil$ and shift each coordinate to
$w_{e,ij}=z_{e,ij}+L_*\in\{0,\ldots,2L_*\}$.
Since every optional base has $q'$ elements, the shifted and signed summed
profiles determine each other by subtracting $q'L_*\mathbf1$.

For the full-row-rank integer representation define
\begin{equation}
 g(y)=\det\!\left(A_c\diag_e(y^{w_e})A_c^\top\right)
 =\sum_{B'\text{ base}}\det((A_c)_{B'})^2
                       y^{\sum_{e\in B'}w_e}.
 \label{eq:matroid-profile-polynomial}
\end{equation}
This is the Binet--Cauchy identity used by
\citet[Lemmas 4.3--4.4]{Berstein2008}. To see it directly, expand
$\det(CD^\top)$ over $q'$-column subsets with
$C=A_c\diag(y^{w_e})$ and $D=A_c$. A dependent subset has zero
minor; a base contributes its squared determinant times its monomial.
Every coefficient is nonnegative and is positive exactly when its profile
is attained. Characteristic-zero arithmetic prevents cancellation. A nonzero residue of an integer coefficient certifies attainability, but a
zero residue need not mean a zero coefficient: positive coefficients can
vanish modulo a prime. An exact zero test therefore needs rational/integer
arithmetic or additional coefficient bounds and sufficient modular recovery.

Each coordinate degree is at most $D_*=2q'L_*$. Evaluate the determinant
on the product grid $\{1,\ldots,D_*+1\}^d$. Univariate polynomials of
degree at most $D_*$ are determined by their values at $D_*+1$ distinct
points; applying this fact coordinate by coordinate makes the tensor
Vandermonde system invertible. Exact rational solution recovers all
coefficients. This is an equivalent realization of the established
interpolation algorithm, whose original presentation uses a moment-curve
substitution. It enumerates at most $R_*=(D_*+1)^d$ positive profiles.

For each positive target profile recover one actual base by deletion.
Start with all optional columns. Test removal of each column by asking
whether that target coefficient in \eqref{eq:matroid-profile-polynomial}
remains positive; delete the column if it does. Keep the original $q'$ rows
in every test, even after a rank drop, which then correctly produces the
zero polynomial. Do not replace the representation by a smaller row basis.
The target stays feasible throughout. A failed deletion cannot become
feasible after further deletions. If the surviving set contained an extra
column outside one target-profile base, that column could be deleted while
retaining the target, contradicting its failed deletion. Consequently the
surviving columns are exactly a base with the target profile. At most $m$
tests suffice. Lift it by \eqref{eq:matroid-contraction} and return it.

\subsection{Coverage, zero information, and bit bounds}
Fix an arbitrary target base $B$ with information rank $r>0$. Its
maximum-volume trial from Lemma~\ref{lem:spectral-normalization} retains
all of $B$. Its forced owners form an independent subset of $B$, and the
retained set has rank $q$ because it contains $B$. Thus the trial survives.
If $q'=0$ it returns $B=F$. Otherwise $B\setminus F$ has a detected profile;
let $\widehat B$ be the recovered lifted base. Equal signed sums imply that
each transformed information entry differs by less than $q'h\le qh$,
because each optional rounding residual is in $[0,q'h)$ and the exact
forced/prior contributions cancel. Their symmetric difference has norm
at most $rqh=\eta$. The forced floor and rectangular reconstruction then
give \eqref{eq:spectral-cover}. They also prove the exact range equality.
For information rank zero, if $Q_0=0$, restrict to exactly zero-matrix
elements and return a base of that restriction only if its rank is still
$q$. This covers all zero-information bases; a nonzero prior excludes them.

There are at most $\binom Mr$ rank-$r$ trials and
\[
 D_*\le2q\lceil4prq/\eta\rceil,
 \qquad
 |\mathcal C|\le1+\sum_{r=1}^p\binom Mr
 \bigl[2q\lceil4prq/\eta\rceil+1\bigr]^{r(r+1)/2}.
\]
For $q'=0$ a trial returns one base, also within this bound. The initial one
permits the zero-information representative. Normalization and contraction
are polynomial rational operations. At a grid point the monomial has bit
length $O(dL_*\log(D_*+1))$. Matrix products and determinants have polynomial
bit length, even for growing $q'$: determinant numerator lengths are bounded
by the sum of row entry lengths plus $O(q'\log q')$. The tensor Vandermonde
system has polynomial dimension for fixed $d$ and polynomial-bit rational
entries, so exact elimination is polynomial. Deletion may conservatively
recompute all coefficients at each of at most $mR_*$ tests; this remains
polynomial. Basis reconstruction and output length are polynomial too.
This completes the proof in the Turing model, with no floating zero test,
independence-oracle conversion, or hidden feasibility relaxation.

\subsection{Exact witnesses for the representation restrictions}
The squared-minor argument is specific to a single represented matroid.
For two representations the mixed determinant contains products of different
minors, whose signs may differ. With row matrices $A_1=(1,1)$ and
$A_2=(1,-1)$, and identical zero profiles, both singleton sets are common
bases but $\det(A_1A_2^\top)=0$. Thus direct substitution loses feasible
profiles. A separate randomized profile algorithm for represented matroid
intersection is established by \citet[Section 3]{Berstein2010}; element
indeterminates and random substitution address that cancellation. We do not
claim its guarantee as part of our deterministic theorem. Likewise,
\[
 \det\begin{pmatrix}1&1&0\\1&0&1\\0&1&1\end{pmatrix}=-2
\]
is nonzero over $\Q$ but zero over $\mathbb F_2$, so reinterpreting a
finite-field representation can change the bases.

Projection does not remove attainability requirements either. The rank-one
uniform matroid on two elements with scalar information $1$ and $3$ has
projected information polytope $[1,3]$. Its integer point $2$ is not attained
by any base. Integer optimization over that interval would therefore solve
a different problem. Finally, the independence-oracle result of
\citet[Theorem 1.1]{Berstein2008} assumes a fixed number of distinct weight
values. The number of our rounded labels grows with input size and
$1/\eta$; substituting them in that oracle theorem does not retain a fixed
polynomial exponent. These examples explain limitations of the proof routes,
not hardness of every possible broader spectral-cover algorithm.
